\section{Ordinary arithmetic on hereditary Cantor forms}\label{sec:ordinary}
Write $W(e)=[e]$, $0=[]$, $1=W(0)$, $\omega=W(1)$.
A natural number $n$ in an ordinal expression denotes the list of $n$ zeros.
The length $|x|$ is its number of outer entries, and $h(x)$ is its first
entry when $x>0$. All variables for ordinals in this section range over
valid finite notations. Thus, for example, a formula about all $e$ is an
arithmetical assertion about codes with a decidable validity guard.
Juxtaposition of ordinal variables denotes multiplication once it has been
defined. We first use concatenation of \emph{lists}, written $U\mathbin{\|}V$.

\subsection{Addition, subtraction on the left, and absorption}
Given $a,b$, put $a+0=a$. For $b>0$ retain from $a$ its longest prefix
with entries $\ge h(b)$ and concatenate $b$. This is the normal form of
ordinal addition. To see what happens with many summands, allow an arbitrary
finite word of exponents and scan it from left to right. Each input letter
$s$ removes the current final letters $<s$, then joins the end of the word.
Write $N$ for the result. In any surrounding word, replacing adjacent letters
$r,s$ with $s$ when $r<s$ preserves $N$. Indeed, inserting $r$ can only
have removed letters smaller than $r$. The subsequent insertion of $s$
would remove all of them, as well as $r$. The normalized stack after $s$
is identical in the two scans. The scan itself implements just such
deletions. It follows that $N(U\mathbin{\|}V)$ is unchanged by normalizing
either factor first. For descending words it is exactly the displayed
prefix rule.

\begin{theorem}\label{th:ordinary-add}
Addition is associative and has two-sided unit $0$. For fixed $a$, the
map $b\mapsto a+b$ is an increasing bijection onto the terms $\ge a$,
with a primitive recursive inverse on its range. For fixed $b$ the map
$a\mapsto a+b$ is weakly increasing. In particular, left cancellation
holds, $a+(b+1)=(a+b)+1$, and $a+1$ appends one zero.
\end{theorem}
\begin{proof}
Associativity follows by normalizing the concatenation of three words,
and the empty word gives the unit. When a positive summand $d$ removes
a suffix of $x$, its first entry exceeds the first removed entry.
When it removes nothing, $x$ is a proper prefix of $x+d$. Thus $x<x+d$.

For $b<c$ discard their longest common prefix and take $d$ to be the
remaining suffix of $c$. If $b$ ends at the common prefix, $b+d=c$ by
concatenation. Otherwise its first remaining entry is smaller than the
first of $d$. Every remaining entry of $b$ is removed and the common
prefix is retained. Again $b+d=c$. For $b=c$ take $d=0$. Hence
$a+c=(a+b)+d>a+b$ whenever $b<c$. This gives uniqueness of the residual
and the stated range.

For weak monotonicity in the first argument it suffices to insert one
letter $s$ at a time. Deleting the entries below $s$ from two ordered
words preserves their lexicographic order because a first mismatch either
survives, becomes a prefix comparison, or disappears in both.
Appending $s$ to both truncated words preserves that order. Indeed, in
the proper-prefix case the next letter of the longer word is at least
$s$. Equality of this next letter with $s$ leaves the shorter result a
proper prefix. An induction on the scanned length proves the assertion.
\end{proof}

\begin{theorem}[Additive indecomposability]\label{th:add-indecomp}
For $x>0$, the singletons $W(e)$ are exactly the terms closed below themselves under sums of two smaller terms. Equivalently, $x$ is a singleton if and only if it cannot be expressed as a sum of two smaller terms. The map $e\mapsto W(e)$ gives
an increasing enumeration of exactly these terms. Also $u+W(e)=W(e)$
whenever $u<W(e)$.
\end{theorem}
\begin{proof}
A nonempty word is below $[e]$ exactly when its first entry is below
$e$. The sum of two such words still uses only entries below $e$, and
insertion of $e$ deletes all of the first word. The empty-word cases follow
from the unit. If $x$ has at least two entries, split off its final
singleton. Both pieces are smaller than $x$ and concatenation is their
ordinal sum. Singleton comparison compares their sole entries.
\end{proof}

\subsection{Multiplication by expanding the right operand}
Set $0b=a0=0$. For $a>0$ replace each entry $e$ of $b$ by the block
\[
 B_a(0)=a,\qquad B_a(e)=W\bigl(h(a)+e\bigr)\quad(e>0),
\]
and add the blocks in their original order. In particular $aW(e)=B_a(e)$.
For $r<s$, all exponents of $B_a(r)$ are smaller than the sole exponent
of $B_a(s)$. For $r=0$ this uses $h(a)<h(a)+s$. For $r>0$ it uses strict
increase of $h(a)+{}\cdot$. Consequently a deletion $r,s\rightsquigarrow s$
in a word of exponents preserves the sum of its expanded blocks.
The first block has first exponent at least as large as every later
block exponent, so it survives all later insertions. Thus
\[
                  h(ab)=h(a)+h(b)\qquad(a,b>0).                 \tag{1}\label{eq:arith-head}
\]

\begin{theorem}\label{th:ordinary-mul}
Multiplication is associative, has unit $1$ and absorbing element $0$,
and satisfies $a(b+c)=ab+ac$. For $a>0$, $b\mapsto ab$ is strictly
increasing and left cancellation holds. Multiplication is weakly increasing
in its first argument. In particular $a(b+1)=ab+a$.
\end{theorem}
\begin{proof}
Expansion of a concatenation gives the sum of its two expansions.
The preceding deletion argument allows the concatenation to be normalized
first, proving the distributive law. Both unit identities follow from
$B_a(0)=a$ and $B_1(e)=W(e)$ for every $e$. Associativity reduces by
finite distributivity to $a(bW(e))=(ab)W(e)$ for each entry $e$.
Zero factors are immediate and $e=0$ is the unit identity.
For $a,b,e>0$ the two sides are monomials with exponents
$h(a)+(h(b)+e)$ and $(h(a)+h(b))+e$, by (1). They agree by associativity
of addition. If $b<c$, take $d>0$ with $b+d=c$. For $a>0$, expansion
gives $ad>0$, so $ac=ab+ad>ab$. If $0<a\le a'$, then $h(a)\le h(a')$. Each expanded block for $a$ is bounded by its counterpart for $a'$.
Addition is weakly increasing in both arguments, so the same holds for
the sum of all blocks. The case $a=0$ follows from leastness of zero.
\end{proof}

For $s,t\ge0$ this calculation gives
$W(s)W(t)=W(s+t)$, including $t=0$. For a positive $a$ it gives
$a\omega=W(h(a)+1)$. If the first exponent $g$ of $a$ has multiplicity
$j$, then $an$, for a positive finite $n$, has $j(n-1)$ copies of $g$
followed by $a$. Each new copy of $a$ deletes exactly the part of the
preceding copy below $g$. This proves the formula by ordinary induction.
In particular, finite ordinals have their usual addition and multiplication.

\begin{corollary}[Multiplicative indecomposability]
The nonzero terms $x$ satisfying
$u,v<x\Rightarrow uv<x$ are exactly $1$, $2$ and the $W(W(e))$.
\end{corollary}
\begin{proof}
The finite cases $1,2$ can be checked directly. For finite $n\ge3$,
$(n-1)^2\ge n$. Below $W(W(e))$ the heads of nonzero factors are
$<W(e)$, and their sum remains $<W(e)$ by Theorem~\ref{th:add-indecomp}.
Equation (1) bounds their product. Conversely an infinite $x$ with at
least two entries and head $g$ strictly exceeds $W(g)$. Both $W(g)$
and the natural number $|x|+1$ are below $x$, but their product, a list
of $|x|+1$ copies of $g$, exceeds $x$. For infinite $x=W(g)$, if
$g$ is not a singleton then $g>0$ splits as $p+q$ with $p,q<g$.
The two smaller monomials $W(p),W(q)$ multiply to $x$.
\end{proof}

\subsection{Exponentiation with arbitrary base}
The following decomposition is also an indexing theorem for the limits.
Every term $b$ has a unique expression
\[
                         b=\omega Q+n,\qquad n<\omega.            \tag{2}\label{eq:arith-omega-div}
\]
Count its final zeros to obtain $n$. For each remaining entry $e>0$,
solve $1+d=e$ by subtracting one if $e$ is finite, and by taking $d=e$
if $e$ is infinite. Indeed, in the latter case the head of $e$ is
positive, so $1$ is absorbed. The solution is unique and ordered as
$e$ varies by Theorem~\ref{th:ordinary-add}. These solutions in their
given order form $Q$. The rules for multiplication replace its entries
$d$ by $1+d$, also when $d=0$. The recovered positive entries of $b$
and its $n$ zeros prove (2) and uniqueness. Thus $Q\mapsto\omega Q$ is
an order isomorphism onto zero together with the nonzero limits (the
nonempty terms with no final zero).

Write $P_a(n)$ for the ordinary finite product of $n$ copies of $a$,
with $P_a(0)=1$. Set $a^0=1$, $0^b=0$ for $b>0$, and $1^b=1$.
For $a\ge2$ define the code $k_a$ and the general power by
\[
 k_a=\begin{cases}1,&a<\omega,\\ h(a)\omega,&a\ge\omega,\end{cases}
 \qquad
 a^{\omega Q+n}=W(k_a Q)\,P_a(n).                              \tag{3}\label{eq:arith-power}
\]
These operations are primitive recursive. All word scans, prefix cuts and
extractions in (1)--(3) are bounded by the input word length. The recursion
$P_a(n+1)=P_a(n)a$ uses the already defined multiplication. A changing
accumulator is allowed in ordinary primitive recursion. All further scans
are on outputs of earlier primitive recursive operations. For instance,
when $a\ge\omega$, $h(a)>0$ and $k_a=W(h(h(a))+1)$ is a single monomial.
Formula (3) therefore also covers bases and exponents of arbitrarily
large finite nesting depth.

\begin{theorem}[Power laws]\label{th:ordinary-pow}
For all terms $a,b,c$,
\[
 a^{b+c}=a^b a^c,\qquad (a^b)^c=a^{bc},\qquad
 a^{b+1}=a^b a,\qquad \omega^b=W(b).
\]
For $a\ge2$ the exponential is strictly increasing in $b$. For fixed
$b$ it is weakly increasing in $a$.
\end{theorem}
\begin{proof}
First fix $a\ge2$ and put $k=k_a$. By (1) and finite induction,
$h(P_a(m))=h(a)m$, including $m=0$. This head is $<k$. For a finite
base it is zero. For an infinite base and $m>0$ it has head $h(h(a))$, below
the exponent of $k$. In consequence, for $C>0$,
\[
             P_a(m)W(kC)=W(kC).                                \tag{4}
\]
Indeed $kC\ge k$, since multiplication by positive $k$ is increasing.
As $k$ is a single monomial, every exponent of the word $h(P_a(m))$
is smaller than its exponent. More explicitly, if $h(P_a(m))>0$, then
$h(h(P_a(m)))<h(k)\le h(kC)$.
Thus $h(P_a(m))+kC=kC$, proving (4) by the block rule.

Write $b=\omega B+m$, $c=\omega C+n$. If $C=0$, appending $c$ changes
the finite tail of $b$ to $m+n$ and $P_a(m+n)=P_a(m)P_a(n)$ by finite
induction. If $C>0$, $m+\omega C=\omega C$ by absorption, so the sum
has quotient $B+C$ and tail $n$. Removing the middle factors by (4)
in $a^b a^c$ leaves
$W(kB)W(kC)P_a(n)=W(k(B+C))P_a(n)$. This is the first power law.
It includes the successor identity because $a^1=a$.

For the second power law take $b>0$ and define
\[
 D=1\quad(b<\omega),\qquad D=B\omega\quad(b=\omega B+m,\ B>0).
\]
We claim $b\omega=\omega D$ and $k_{a^b}=kD$. Here $a^b\ge2$ because positive
finite powers of $a$ are at least $a$, and $B>0$ gives first exponent
$\ge kB>0$ by (1). If $b$ is finite then $b\omega=\omega$.
A positive finite power of a finite base is finite. For an infinite base,
$h(P_a(b))=h(a)b$ has the same head as $h(a)$, so
$h(P_a(b))\omega=h(a)\omega$. This proves both claims in the finite case.
If $B>0$, equation (1) gives $h(b)=1+h(B)$, since adding the finite
tail preserves the head. Now $b\omega=W((1+h(B))+1)=\omega(B\omega)$.
The head $H=h(a^b)=kB+h(P_a(m))$ has the same head as $kB$ because the second
summand is $<k$, whereas $kB\ge k$, with $k$ a monomial.
Thus $H\omega=(kB)\omega=kD$, proving the second claim.

Apply the first power law repeatedly to obtain
$P_{a^b}(n)=a^{bn}$ by finite induction on $n$.
Since $bc=b\omega C+bn=\omega(DC)+bn$, the first law and (3) give
\[
       a^{bc}=W(kDC)P_{a^b}(n)=(a^b)^c .
\]
If $b=0$, both sides of the second law are $1$. For bases $0$ and $1$
the laws follow by splitting zero from positive exponents and using that
a sum is zero exactly when both summands are zero, and a product is zero
exactly when a factor is zero.

For $a=\omega$, $k=\omega$ and $P_\omega(n)=W(n)$ by the monomial
product formula. Thus (3) is $W(\omega Q)W(n)=W(b)$. For $b<c$, write
$c=b+d$ with $d>0$. Since $a^d>1$ and $a^b>0$ for $a\ge2$, strict
increase of multiplication in the right argument gives $a^c>a^b$.
Finally $k_a$ is weakly increasing for $a\ge2$, by monotonicity of the
head and of multiplication, and $P_a(n)$ is weakly increasing in $a$
by finite induction. Formula (3) proves weak monotonicity in the base.
The exceptional bases satisfy it directly.
\end{proof}

\subsection{Continuity with explicit witnesses}
For a strictly increasing function $F$ on the codes and a nonzero limit
$L$, the limit equation $F(L)=\sup_{d<L}F(d)$ has the following
witnessed meaning. Every $F(d)$ with $d<L$ lies below $F(L)$, and for
every $z<F(L)$ we compute $d<L$ with $z<F(d)$. These two assertions
imply the least-upper-bound property by the decidable linear order.

\begin{theorem}\label{th:ordinary-cont}
The maps $d\mapsto a+d$, $d\mapsto ad$ for $a>0$, and $d\mapsto a^d$
for $a\ge2$ satisfy these limit equations at every coded nonzero limit.
\end{theorem}
\begin{proof}
Strict increase proves the first assertion. For addition, given
$z<a+L$, use $d=0$ if $z<a$. Otherwise compute $q$ with $a+q=z$. Then $q<L$, and $d=q+1<L$ works since $L$ is a limit.
For multiplication, use the Euclidean division theorem for the codes
proved in this paper. Write $z=aq+r$, $r<a$. Then $aq\le z<aL$ gives
$q<L$, and $z<a(q+1)$ with $q+1<L$.

For powers, (2) writes the limit as $L=\omega Q$ with $Q>0$. Thus
$a^L=W(kQ)$ where $k=k_a$. If $z=0$, use $d=0$. If $z>0$, let
$t,r$ be the Euclidean quotient and remainder of $h(z)$ on division
on the left by $k$. Then $h(z)=kt+r<kQ$, $r<k$, and $t<Q$.
If $a$ is infinite, write $g=h(h(a))$, so $k=W(g+1)$. Put
$n=|r|+1$. Every exponent of $r$ is $\le g$, and $h(a)n$ begins with
at least $n$ copies of $g$. Thus $h(a)n>r$. For $d=\omega t+n$
equation (1) now gives
\[
 h(a^d)=kt+h(P_a(n))=kt+h(a)n>kt+r=h(z).
\]
If $2\le a<\omega$, then $k=1$, $r=0$, and $t=h(z)$. Use
$n=|z|+1$. The word $a^{\omega t+n}=W(t)P_a(n)$ consists of
at least $n$ copies of $t$, since the natural number $a^n\ge n$.
It exceeds $z$ by length or at the first smaller exponent of $z$.
In either case $\omega t+n<\omega(t+1)\le\omega Q$. This follows
from $n<\omega$ and left distributivity. All three witnesses have
been constructed by primitive recursive operations.
\end{proof}

For $a=1$ the power function is constant $1$, and for base $0$ the
positive powers are constant $0$. Including the exponent $0$ makes the
supremum for this base equal to $1$, since $0^0=1$. Taking the supremum over
$0<d<L$ gives the article's limit equation in these cases too. The zero
multiplication map has constant value $0$ and also preserves these suprema.

\paragraph{Calculations.}
We have $1+\omega=\omega<\omega+1$ and
$2\omega=\omega<\omega 2$. Thus right cancellation for addition and
multiplication can fail. Also $(1+1)\omega=\omega$ while
$1\omega+1\omega=\omega 2$, exhibiting failure of the other distributive
law. Formula (3) yields $2^\omega=\omega$, $2^{\omega+1}=\omega 2$,
and $\omega^2=W(2)>\omega$. In particular, exponentiation also
fails to commute. With the left argument varying up to $\omega$, the
values $n+1$, $n2$ and $n^2$ on finite arguments are cofinal in
$\omega$, whereas $\omega+1$, $\omega 2$ and $\omega^2$ all exceed it.
This proves failure of continuity in the other argument in all three cases.

The spiral in the opening illustration of \emph{Ordinal number} uses
$f(x)=\omega(1+x)$. Its least fixed point is $W(\omega)=\omega^\omega$.
Indeed, $f(x)>x$ for finite $x$. For $\omega\le x<W(\omega)$ one has
$1+x=x$ and the positive finite head $h(x)$ increases to $1+h(x)$
under multiplication by $\omega$ on the left. At $x=W(\omega)$ the
head becomes $1+\omega=\omega$ and the product is exactly $W(\omega)$.

\subsection{Prime factors and a unique ordered factorization}
An ordinal prime is a term $x>1$ with no expression $x=ab$ with $a,b<x$.
This condition concerns equality. Closure of an initial segment under
all products, treated above, is stronger.

\begin{theorem}[Prime classification and normal factorization]\label{th:ord-primes}
The primes of this notation system are the finite primes, the limit primes
$W(W(d))$, and the infinite successor primes $W(d)+1$ with $d>0$.
Every positive term has a unique prime factorization subject to the
following ordering rules. All limit factors precede all other factors,
the limit factors are nonincreasing, and each uninterrupted group of
finite primes is nonincreasing. All the factors can be computed
primitive recursively.
\end{theorem}
\begin{proof}
If $ab$ is positive and finite, both factors are finite by (1) or
by monotonicity in each argument. The finite primes are therefore
exactly those of elementary arithmetic. The closure assertion above
already proves primality of $W(W(d))$. Let $x=W(d)+1$ with $d>0$
and suppose $x=ab$ with both factors smaller. Each factor must exceed
$1$. If $b$ is a limit, every block of $ab$ is a singleton with
positive exponent, whereas $x$ ends in zero. If $b$ is finite,
$b\ge2$. Then $a$ must be infinite and the formula for $ab$ has at
least two copies of its leading exponent, whereas $x$ has exactly one.
If $b$ is an infinite successor, $ab$ has at least one positive
singleton block before the last copy of $a$. The singleton survives
because its exponent strictly exceeds $h(a)$. As $a>1$ and $ab$ ends in
zero, $a$ ends in zero too and has at least two entries. Thus $ab$ has
at least three entries, whereas $x$ has two.

We now give a factorization algorithm. Combine adjacent equal entries
of $x$ into runs with exponents $e_0<\cdots<e_s$ listed from the
\emph{right} and multiplicities $m_0,\ldots,m_s>0$. For each $i>0$
compute $d_i>0$ with $e_{i-1}+d_i=e_i$ by the prefix algorithm. Write
the word of $e_0$ as $[g_1,\ldots,g_j]$. The ordered factorization,
temporarily keeping finite factors unexpanded, is
\[
 \left(\prod_{t=1}^{j}W(W(g_t))\right)m_0
       \prod_{i=1}^{s}\bigl((W(d_i)+1)m_i\bigr),               \tag{5}
\]
with products read from left to right. The first product is $W(e_0)$
by the monomial product formula and equals $1$ if $e_0=0$.
Multiplication by $m_0$ produces the bottom run. If the portion already
formed is $v$ with head $e_{i-1}$, then
$v(W(d_i)+1)=W(e_i)+v$. Its leading exponent occurs once, so multiplying
by $m_i$ yields exactly $m_i$ copies of $e_i$ followed by $v$.
Finite induction proves (5). Replace each $m_i$ by its finite prime
factors in nonincreasing order, omitting $1$.

Conversely, any product obeying the ordering rules has an initial
block of limit primes, followed by finite blocks separated by the
successor factors $W(d_i)+1$. Finite blocks are allowed to be empty.
Multiplying exactly as above recovers the run exponents and
multiplicities. The bottom exponent records the initial block
letter for letter. Each next exponent determines $d_i$ uniquely by
left cancellation. Each multiplicity recovers its finite block by
unique natural factorization. This proves uniqueness even when a
multiplicity is one.

To see that no further primes occur, take a product (5) containing at
least two prime factors. Its proper prefix before the last factor is
smaller than $x$ by strict increase of multiplication in the right
argument. Its last prime is also smaller. For a finite last prime this
is immediate if $x$ is infinite, and is ordinary arithmetic otherwise.
If the last factor is a limit, the entire list consists of at least two
nonincreasing limit primes, so the bottom exponent $e_0$ exceeds its
last singleton $W(g_j)$. If the last factor is $W(d)+1$, the preceding
product $v$ exceeds $1$ and $x=W(h(v)+d)+v$. Here $h(v)+d\ge d$. Even in the case of equal leading exponents, the tail $v$ exceeds $1$.
Thus (5) splits into two strictly smaller factors. A single prime
factor and the empty product account for all remaining cases.

For the elementary arithmetic used here, factor a natural $m>1$ by
trial division up to $m$, removing a least divisor greater than one
at each stage. Each such divisor is prime by minimality, and at most
$m$ nonterminal divisions occur since the remaining quotient decreases.
Uniqueness uses Euclid's lemma. If a prime $p$ does not divide $v$, the
ordinary Euclidean algorithm on $(p,v)$ returns gcd $1$, since the only
positive divisors of $p$ are $1,p$. Back-substitution gives integers
$u,w$ with $up+wv=1$. Initialize the two inputs by their two unit
coefficient vectors and replace a row by the difference of that row
and the quotient times the other. Multiplication by $z$ shows that
$p\mid vz$ implies $p\mid z$. Euclid's lemma locates an equal prime
in any second factor list, after which cancellation and finite induction
prove uniqueness. The decreasing Euclidean remainders permit the fixed
fuel $p+v+2$. All scans, products, and sorting steps in (5) are bounded
iterations on finite codes.
\end{proof}
